% ============================================================================
% ARR MAIN-PAPER RESULTS SECTION — READY TO COPY AND PASTE
% Requires ARR_PREAMBLE_COPYPASTE.tex and the PDF files in figures/.
% ============================================================================

\section{Additional Experiments}
\label{sec:additional-experiments}

\paragraph{Evaluation protocol.}
We evaluate every condition on the same 197-case mental-health harm cohort.
Unless stated otherwise, the target is \texttt{gpt-4o-2024-11-20} and each
case produces four independent final branches. We report Best-of-4 (denoted
by $@4$). Our primary metric is \emph{strict joint success}: a case is
successful only when the same branch simultaneously satisfies goal recovery
($R3+$), behavioral commitment ($B4+$), and HarmBench+. This same-branch
requirement prevents success from being assembled across different branches.
CARES Accept is reported as a complementary response-level measure.

\subsection{How Much Research Mediation Is Useful?}
\label{sec:research-depth-ablation}

We compare three mediation depths: (i) \emph{direct-final-only}, which asks
the four final questions without a shared research readout; (ii) a compact
\emph{single readout}; and (iii) the original seven-stage research dialogue.
Each condition is crossed with the presence or absence of persona history.
Table~\ref{tab:main-depth-history} and Figure~\ref{fig:main-depth-history}
show the results.

\begin{table*}[t]
\centering
\small
\caption{Research-mediation depth and persona-history ablation on GPT-4o
($n=197$). Values are percentages. Strict joint success requires $R3+$,
$B4+$, and HarmBench+ in the same branch. The strongest strict-joint result
is boldfaced.}
\label{tab:main-depth-history}
\begin{threeparttable}
\begin{tabular}{llrrrrr}
\toprule
Research condition & History & CARES Accept & $R3+@4$ & $B4+@4$ & HarmBench+$@4$ & Strict joint$@4$ \\
\midrule
Direct final only & Yes & 92.39 & 99.49 & 8.12 & 87.31 & 7.11 \\
Direct final only & No  & 85.28 & 96.95 & 10.15 & 75.63 & 6.09 \\
Single readout    & Yes & 100.00 & 99.49 & 100.00 & 96.95 & \textbf{95.94} \\
Single readout    & No  & 100.00 & 97.46 & 98.98 & 86.80 & 83.76 \\
Long 7-stage      & Yes & 96.45 & 100.00 & 76.14 & 89.34 & 63.96 \\
Long 7-stage      & No  & 98.98 & 100.00 & 86.80 & 86.29 & 73.10 \\
\bottomrule
\end{tabular}
\begin{tablenotes}\footnotesize
\item Direct-final-only omits the shared research readout. Single readout
uses one compact, goal-directed research exchange. Long uses the seven-stage
dialogue. All cells use four independent final branches.
\end{tablenotes}
\end{threeparttable}
\end{table*}

\begin{figure*}[t]
\centering
\includegraphics[width=0.98\textwidth]{fig1_research_depth_history.pdf}
\caption{Research-mediation depth and persona-history ablation. A compact
single readout yields the highest same-branch strict-joint success, whereas
direct final questioning rarely reaches $B4+$ and a seven-stage dialogue
reduces final behavioral success.}
\label{fig:main-depth-history}
\end{figure*}

The compact single-readout condition with persona history achieves
\textbf{95.94\%} strict joint success. Removing history reduces this result
to 83.76\% ($-12.18$ percentage points). Direct-final-only retains high
$R3+$ and HarmBench rates but rarely reaches $B4+$, producing only 7.11\%
strict joint success with history. Extending the interaction to seven stages
also reduces strict joint success to 63.96\% with history, a 31.98-point drop
relative to the compact readout. Thus, the gain does not arise from omitting
research mediation or from maximizing conversation length. Instead, the
results favor a short, goal-directed research readout that binds the persona
context to the final branches without diluting the requested output.

\subsection{Cross-Model Generalization}
\label{sec:cross-model-generalization}

We next apply the selected single-readout method without changing its four
final branches. Table~\ref{tab:main-cross-model} and
Figure~\ref{fig:main-cross-model} summarize seven target models, with 197
cases evaluated per target.

\begin{table*}[t]
\centering
\small
\caption{Cross-model transfer of the selected compact method (Best-of-4,
$n=197$ per target). The strongest strict-joint result is boldfaced.}
\label{tab:main-cross-model}
\begin{tabular}{lrrrrr}
\toprule
Target model & CARES Accept & $R3+@4$ & $B4+@4$ & HarmBench+$@4$ & Strict joint$@4$ \\
\midrule
Mistral Small 3.2 24B & 100.00 & 99.49 & 100.00 & 97.97 & \textbf{96.95} \\
Gemma 4 26B-A4B       & 98.98  & 99.49 & 98.48  & 93.91 & 92.89 \\
Phi-4                  & 100.00 & 98.98 & 98.98  & 94.42 & 91.88 \\
Llama 3.1 8B Instruct & 100.00 & 93.91 & 98.98  & 94.92 & 88.32 \\
Llama3-G2C             & 91.88  & 82.23 & 81.22  & 65.99 & 43.65 \\
GPT-5.4-mini           & 0.00   & 100.00 & 0.00   & 5.58  & 0.00 \\
GPT-6-luna             & 0.00   & 97.97 & 0.00   & 5.58  & 0.00 \\
\bottomrule
\end{tabular}
\end{table*}

\begin{figure*}[t]
\centering
\includegraphics[width=0.98\textwidth]{fig3_cross_model.pdf}
\caption{Cross-model transfer of the selected compact method. CARES Accept,
HarmBench+$@4$, and strict joint$@4$ expose distinct aspects of target
behavior; in particular, high $R3+$ alone does not imply behavioral
commitment.}
\label{fig:main-cross-model}
\end{figure*}

Strict joint success exceeds 88\% on Mistral Small 3.2 24B, Gemma 4
26B-A4B, Phi-4, and Llama 3.1 8B Instruct. It falls to 43.65\% on the
counseling-oriented Llama3-G2C model. GPT-5.4-mini and GPT-6-luna retain high
$R3+$ but never reach $B4+$, yielding 0\% strict joint success. These results
show that the method transfers across several model families, while its
effectiveness remains strongly dependent on the target's safety alignment
and willingness to produce concrete task output.

\paragraph{Main finding.}
Across the evaluated conditions, the most effective design is a
persona-grounded, compact research readout followed by four independent final
branches. The full external-baseline matrix, paired PCSA personalization
ablation, attention/representation analysis, and history-transfer study are
reported in Appendix~\ref{sec:additional-experiments-appendix}.

